\chapter{Problem Statement}

\section{Rationale}
The public transportation system has long been a commodity that comes in the form of bus transport, train networks, and the like, with predetermined routes and fares \cite{Tironi2018}. In Metro Manila, traffic congestion has long been a problem stemming from insufficient road infrastructure and poor traffic mitigation policies, resulting in an inefficient public transport system \cite{Spatio2026}. One of the government's actions for an alternative transportation system that runs within the metropolitan area is the implementation of the EDSA Carousel in 2020. It is the only busway system in Metro Manila implemented by the government to have a designated lane, separated from other vehicles by concrete barriers \cite{Chua2026}. Over the years, the number of commuters using the EDSA Carousel has steadily risen, averaging approximately $1.83\times10^5$ daily passengers, with single-day peaks exceeding $3.2\times10^5$ during high-demand periods, reaching a new high of nearly 67 million in 2025, up from 63.02 million in 2024 
\cite{DOTr2025Ridership}. 


However, despite its growing demand, the system still faces several problems that can affect the quality of bus routing. To start, the LTFRB relied on a tedious manual monitoring process by deploying personnel at each stop \cite{EDSApolicy2023}. There are also certain portions where it operates with mixed traffic, causing delays among buses. It is also clear that bus bunching has not been resolved; most of the time, their designated dispatch schedule is not followed \cite{Ollero2024EDSA}. Moreover, the operation of the bus system is also largely influenced by external factors, such as weather and road maintenance. Sections of the EDSA Bus Carousel routes typically close during tropical cyclones to prevent potential damage from continuous rain and strong winds \cite{DOTr2020Suspension}. Given these operational issues, the bus dispatching process of the EDSA Bus Carousel can be modeled as a Vehicle Scheduling Problem (VSP), in which buses must be dispatched and coordinated efficiently despite these issues. By improving the scheduling, operational issues, such as bus bunching, irregular dispatch intervals, and passenger waiting times, will be reduced, thereby improving the overall bus dispatch efficiency.

Busway systems operate in a highly dynamic, unpredictable environment. In the past, Vehicle Scheduling Problems (VSP) relied on traditional scheduling techniques that focused on finding the optimal path by determining the shortest distance. Static scheduling is the pre-computed timetabling that fixes departure times before operation begins, with no mechanism to revise them in response to real-time conditions \cite{Ju2023Joint,Cao2022Train,Nie2025CMRM}. However, because transit environments are inherently stochastic, these static schedules perform poorly during daily operations, which causes severe operational issues such as bus bunching, node failures, and an increase in passenger wait times \cite{Daganzo2009,FEMbusbunching}.

To address the rigidity of static schedules, recent scheduling studies have shifted toward dynamic, AI-driven approaches, with Multi-Agent Reinforcement Learning (MARL) as the leading solution. Past studies have successfully utilized MARL to overcome the scalability limits of earlier systems by distributing the scheduling decisions across autonomous agents, improving the overall system's adaptability to real-time traffic \cite{Rodriguez2023Cooperative,Wang2023MultiObj,Wangsun}. However, a significant gap remains in how these state-of-the-art models are evaluated. Existing literature primarily validates the MARL scheduling models in a controlled, simulated environment. Previous evaluations rarely account for the heavy-tailed, asymmetric disruptions characteristic of actual transit operations, leaving their performance under non-ideal, real-world traffic anomalies largely untested. To the best of our knowledge, no existing MARL bus-scheduling model has been sufficiently tested against severe, simultaneous disruptions such as sudden passenger surges and weather-induced delays. To bridge this gap, this study develops and evaluates a MARL-based bus scheduling model within a stochastic traffic simulation subjected to non-ideal conditions.

This study develops and evaluates a MARL-based bus scheduling model on a calibrated stochastic simulation of the EDSA Carousel corridor, with non-ideal conditions operationalized as the simultaneous and independent injection of passenger demand surges, severe-weather travel-time skew, and Poisson-distributed bus breakdowns. The objective is to quantify how MARL performance degrades from its ideal-condition baseline as disturbance intensity rises, producing the first integrated robustness characterization for MARL bus scheduling under this specific combination of disturbances.

\section{Research Gap}
Existing Multi-Agent Reinforcement Learning frameworks for bus scheduling demonstrate measurable improvements in passenger waiting time and headway regularity, but these gains are validated almost exclusively in simulation environments that assume stationary passenger arrivals, symmetric Gaussian travel times, and a failure-free fleet. The performance of MARL bus-scheduling models has not been characterized when these idealizing assumptions are simultaneously violated by (i) stochastic passenger demand surges, (ii) heavy-tailed weather-induced travel-time delays, and (iii) discrete bus breakdowns. Without this characterization, it cannot be determined whether reported MARL gains persist, degrade gracefully, or collapse under realistic operating disturbances, which in turn blocks the transition of MARL bus scheduling from simulation to real urban transit deployment.

The weather-disturbance class (W) in particular was identified through the literature survey conducted earlier in this study (Section~\ref{subsec:marl-applied}), which found that no prior MARL bus-scheduling paper models heavy-tailed weather-induced travel-time delays (Table~\ref{tab:marl_performance}, column W). Its operational relevance to the EDSA corridor is established by the rainfall-driven reductions in average speed and free-flow capacity documented in Section~1.1 \cite{TSSP_Rain2018} and by the typhoon-related service suspensions recorded for the corridor \cite{DOTr2020Suspension}. The lognormal parameterization adopted for this disturbance class follows the Kolmogorov--Smirnov-validated form of Patil et al.~\cite{Patil2025Conformal}, introduced in this study to address the resulting lack of temporally aligned, corridor-specific anomaly data (Section~\ref{subsec:disturbance-gap}).



\section{Research Objectives}
\subsection{Main Objective}
To investigate the performance of a developed Multi-Agent Reinforcement Learning (MARL) scheduling model for dynamic bus scheduling under non-ideal traffic conditions on a calibrated simulation of the EDSA Carousel corridor.

\subsection{Specific Objectives}
\begin{enumerate}
   \item \textbf{Environment Construction:} To construct a stochastic traffic simulation environment representing a defined traffic area.

   \textbf{Expected Output 1.1:} A simulation model of a defined operational scope, validated against empirical traffic data to establish an environment that accurately approximates real-world traffic.

   \textbf{Expected Output 1.2:} An environment simulator capable of introducing environmental variability through adjustable parameters to simulate non-ideal operating conditions.

   \item \textbf{Algorithm Implementation:} To develop and implement a Multi-Agent Reinforcement Learning (MARL) approach for dynamic bus scheduling.

   \textbf{Expected Output 2.1:} A formulated agent architecture defining individual bus units as agents with discrete state and action space with a continuous reward function that punishes irregularities and passenger service delays.

   \item \textbf{Performance Evaluation Under Ideal and Non-Ideal Conditions:} To evaluate the scheduling performance of the MARL approach under both ideal operating conditions and non-ideal conditions.

   \textbf{Expected Output 3.1:} A quantitative and statistical assessment of MARL scheduling performance under ideal conditions, measuring passenger waiting time, travel time, and headway coefficient of variation, tested against baseline methods.

   \textbf{Expected Output 3.2:} A quantitative and statistical assessment of MARL scheduling performance under non-ideal conditions, measuring passenger waiting time, travel time, and headway coefficient of variation, reported as degradation relative to ideal-conditions performance.
\end{enumerate}


\section{Significance of the Study}
This study contributes both practical and scientific significance.
 
\textbf{Practical significance.} The EDSA Carousel carries on the order of $1.8\times10^5$ passengers daily, with single-day peaks exceeding $3.2\times10^5$ \cite{DOTr2025Ridership}. At this scale, even a small per-passenger improvement aggregates into a large operational effect: an illustrative 30-second reduction in mean passenger waiting time, applied across average daily ridership, corresponds to roughly $1.5\times10^3$ passenger-hours saved per day. (This figure is illustrative of the scale of impact rather than a performance target; actual savings depend on how uniformly the reduction is realized across riders and operating periods.) Beyond average-case waiting time, this study reports controller performance as a function of disturbance intensity, characterizing how a learned policy degrades as operating conditions worsen --- a perspective directly relevant to transit operators making service-reliability decisions during disruptions.

\textbf{Scientific significance.} The primary scientific contribution is empirical: this study produces the first characterization of whether reported MARL bus-scheduling gains persist, degrade gracefully, or collapse when the idealizing assumptions of prior work are violated simultaneously by demand surges, traffic-speed perturbations, weather-induced heavy-tailed delays, and discrete breakdowns. To produce this characterization reproducibly, the study contributes a stress-testing protocol for bus-scheduling reinforcement learning: a calibrated corridor model paired with four configurable stochastic anomaly generators. Because the protocol fixes the corridor, the disturbance generators, and the evaluation metrics, subsequent MARL work can be benchmarked on the same basis, allowing performance claims to be compared rather than reported in isolation. The study further reports the performance-degradation curve of MARL as a function of disturbance intensity, characterizing robustness rather than only average-case performance.

\section{Scope and Limitations}
 
\textbf{Scope.} This study develops and evaluates a MARL-based bus scheduling framework for the EDSA Carousel corridor. The framework is built on a calibrated SUMO microsimulation and runs over a single-day operational horizon; the simulation horizon, fleet size, and stop set are specified in Chapter 3. Each bus is modeled as an independent agent sharing one learned policy, observing only local conditions and acting at designated control stops. The MARL policy is benchmarked against three baselines — No Control, Forward Headway, and Even Headway — using passenger waiting time, total travel time, and headway coefficient of variation as metrics, over at least 30 Monte Carlo runs per regime. The same trained policy is then evaluated under two regimes, ideal and non-ideal, where non-ideal conditions are defined below.

For this study, \textbf{non-ideal conditions} are operating regimes in which one or more of four disturbance classes are active: (i) passenger demand surges, where boarding demand is scaled up above its baseline; (ii) traffic-speed perturbations, where corridor cruising speed varies around its baseline to represent everyday congestion friction; (iii) severe-weather travel-time delays, where inter-stop travel times are drawn from a heavy-tailed distribution to capture rain-induced slowdowns; and (iv) discrete bus breakdowns, where individual vehicles fail at random and leave service. The exact distributions and parameter ranges for each class are specified in Chapter~3. The severe-weather generator characterizes controller behavior across a range of travel-time variability magnitudes rather than across named meteorological categories; mapping specific intensity values to operational weather labels (e.g.\ ``moderate rain'' versus ``typhoon-grade'') is deferred to follow-on work that pairs the framework with a corridor-specific severe-weather travel-time dataset.
\clearpage
\textbf{Delimitations.} (a) Due to computational constraints, the simulation is restricted to a defined operational sub-segment of the EDSA Carousel corridor rather than the entire metropolitan road network, and minor feeder roads leading into the corridor are not modeled. Both restrictions are justified by the same structural fact: the EDSA Carousel operates on a physically separated, barrier-protected busway \cite{Chua2026}, so the agents' state and reward depend only on bus dynamics within the dedicated lane, specifically headways, dwell times, and onboard loads, none of which are directly observed by or computed from feeder-road traffic. Feeder roads affect the corridor only indirectly, through the passenger arrival rates they produce at each stop, and that effect is already captured by the calibrated per-stop demand distributions (Section~3.2.5) without needing to simulate the feeder network itself. Modeling feeder roads in SUMO would add computational cost without adding any new information the agents' observation or reward could use, since the sub-corridor restriction also preserves 1:1 empirical traffic volumes for GEH calibration without resorting to flow scaling; corresponding GEH calibration statistics are reported in Chapter~4. (b) MARL control actions are applied only at designated control stops, selected following criteria established in the literature (see Chapter~3). (c) The study is entirely simulation-based; no on-road deployment is performed. (d) Weather and breakdown effects are modeled as statistical distributions of their impact on travel time and fleet availability, not as physical or mechanical simulations of rain dynamics or vehicle subsystems. (e) Only the four disturbance classes defined above are considered; other disturbances (road accidents, signal failures, driver behavior variability) are out of scope. (f) Sim-to-real transfer is out of scope; the experimental claims are bounded to the calibrated SUMO microsimulation.